\chap{A complete regularization diagnostic for the escaping model}\label{sec:diagonal-regularization}
The abstract alternatives become explicit for the diagonal model. They also answer a natural objection to the hard prefix corrections: perhaps a smoother correction could fit increasingly well without escaping. It can fit more efficiently, but it cannot avoid the obstruction.

\subsection{The entire regularized family is explicit}
Use $Tu=(u_n/n)$ and $r=w=(1/n)$ as in Section~\ref{sec:escape}. Choose $\alpha=m^{-2}$ for an integer $m\ge1$. The normal equation in each coordinate gives
\[
 a^{(m)}_n=\frac{m^2}{m^2+n^2},\qquad
 v^{(m)}_n=r_n-(Ta^{(m)})_n=\frac{n}{m^2+n^2}.
\]
The repaired ambient candidate is $(-a^{(m)},1)$. Its protected coordinate is one for every $m$.

\begin{proposition}[Energy, residual, and certified tails]\label{WR10}
The displayed $a^{(m)}$ belongs to $\ell^2$ and is the unique regularized correction at $\alpha=m^{-2}$. Set
\[
 A_m=\norm{a^{(m)}}^2=\sum_{n\ge1}\frac{m^4}{(m^2+n^2)^2},\qquad
 B_m=\norm{v^{(m)}}^2=\sum_{n\ge1}\frac{n^2}{(m^2+n^2)^2}.
\]
Then
\begin{equation}\label{eq:soft-bounds}
 \frac m4\le A_m\le\frac{4m}{3},\qquad
 \frac1{4(m+1)}\le B_m\le\frac2m.
\end{equation}
For the sums $A_{m,K},B_{m,K}$ through any integer $K\ge m$,
\begin{equation}\label{eq:soft-tails}
 0\le A_m-A_{m,K}\le\frac{m^4}{3K^3},\qquad
 0\le B_m-B_{m,K}\le\frac1K.
\end{equation}
\end{proposition}
\begin{proof}
For $n>m$, the correction square is at most $m^4/n^4$, whose sum converges. The residual belongs to $\ell^2$ because $T$ is bounded and $r\in\ell^2$. The coordinate normal equation is
\[
 \frac1n\frac{n}{m^2+n^2}=m^{-2}\frac{m^2}{m^2+n^2}.
\]
It holds in $\ell^2$, so WR06 proves global optimality.

For $n\le m$, the correction coordinate is at least $1/2$, yielding $A_m\ge m/4$. Bound its first $m$ squares by one and use
$\sum_{n>m}n^{-4}\le\int_m^\infty x^{-4}\,dx=1/(3m^3)$ for the tail. This gives $A_m\le m+m/3$.

The residual square on $n\le m$ is at most $n^2/m^4$, whose sum is at most $1/m$. The remaining sum is at most $\sum_{n>m}n^{-2}\le1/m$. For its lower bound, $n>m$ implies $(m^2+n^2)^2\le4n^4$, so
\[
 B_m\ge\frac14\sum_{n>m}\frac1{n^2}\ge\frac1{4(m+1)}.
\]
The same comparisons beyond $K$ give \eqref{eq:soft-tails}. All tail bounds concern the actual infinite sums.
\end{proof}
\written{WR10}

Thus the correction norm grows on the order of $\sqrt m$, the residual norm decays on the order of $m^{-1/2}$, and the penalty $\alpha A_m=A_m/m^2$ tends to zero. The smooth correction avoids the hard discontinuity at a prefix boundary, but it still approaches the forbidden constant sequence coordinate by coordinate.

\fixedplate{35_regularized_escape}{At each fixed coordinate, the soft correction approaches one as $m$ tends to infinity. Its transition moves outward as $m$ increases. The chart compares three exact coordinate formulas $m^2/(m^2+n^2)$, not finite approximations to an unspecified optimization procedure. Every displayed sequence has finite energy, but the family's energies are unbounded.}{fig:regularized-escape}

\subsection{The sharp asymptotic correction cost}
The preceding bounds establish divergence and decay. The leading constant can also be computed, which permits a meaningful comparison with hard prefix repair.

\begin{theorem}[A sharp budget asymptotic for the diagonal datum]\label{WR11}
As $m\to\infty$,
\[
 \frac{A_m}{m}\longrightarrow\frac\pi4,\qquad
 mB_m\longrightarrow\frac\pi4.
\]
For the exact budget profile of this infinite problem,
\begin{equation}\label{eq:sharp-budget}
 \lim_{M\to\infty}M\beta_{T,w}(M)=\frac\pi4.
\end{equation}
The hard prefix correction at budget $M=\sqrt N$ instead has
\[
 \sqrt N\,\norm{w-T\mathbf1_N}\longrightarrow1.
\]
\end{theorem}
\begin{proof}
The two normalized sums are Riemann sums on the positive half line:
\[
 \frac{A_m}{m}=\frac1m\sum_{n\ge1}\frac1{(1+(n/m)^2)^2},\qquad
 mB_m=\frac1m\sum_{n\ge1}\frac{(n/m)^2}{(1+(n/m)^2)^2}.
\]
On every bounded interval, ordinary Riemann convergence applies. Beyond $L>1$, the first function is bounded by $x^{-4}$ and the second by $x^{-2}$. Their sum tails are bounded uniformly for sufficiently large $m$ by constants times $L^{-3}$ and $L^{-1}$, respectively. The corresponding integral tails have the same bounds. Letting $m\to\infty$ first on the bounded interval and then $L\to\infty$ proves convergence of the full sums.

The substitution $x=\tan\theta$ gives
\[
 \int_0^\infty\frac{dx}{(1+x^2)^2}
 =\int_0^{\pi/2}\cos^2\theta\,d\theta=\frac\pi4,
\]
\[
 \int_0^\infty\frac{x^2\,dx}{(1+x^2)^2}
 =\int_0^{\pi/2}\sin^2\theta\,d\theta=\frac\pi4.
\]
Put $M_m=\sqrt{A_m}$. These budgets increase strictly to infinity, and $M_{m+1}/M_m\to1$ by the first asymptotic. WR07 gives $\beta(M_m)=\sqrt{B_m}$, so
$M_m\beta(M_m)=\sqrt{A_mB_m}\to\pi/4$.
For $M_m\le M\le M_{m+1}$, monotonicity of $\beta$ yields
\[
 M_m\beta(M_{m+1})\le M\beta(M)\le M_{m+1}\beta(M_m).
\]
Both bounds tend to $\pi/4$, proving \eqref{eq:sharp-budget}. Finally, the earlier telescoping tail enclosure gives
\[
 \frac{N}{N+1}\le N\sum_{n>N}\frac1{n^2}\le1,
\]
which proves the hard-prefix limit.
\end{proof}
\written{WR11}

At large matched correction budgets, the optimal soft repair therefore has about $\pi/4$ times the hard-prefix residual. This improves the efficiency of approximation, not the existence result. Both residuals tend to zero only as the budget tends to infinity. The limiting constant is a calculation for this declared model, not a universal regularization efficiency factor.

\fixedplate{36_budget_product}{The plotted intervals enclose the actual infinite quantities $\sqrt{A_mB_m}$, using finite sums plus the proved tails. The corresponding hard-prefix intervals enclose $\sqrt{N\sum_{n>N}n^{-2}}$ with $N=m$. These have different natural budget sequences. Their limits give the matched-budget asymptotic constants $\pi/4$ and one.}{fig:budget-product}

\subsection{The same operator has an attainable datum}
Change only the residual to $b_n=1/n^2$. Its minimum norm preimage is $a^\dagger_n=1/n\in\ell^2$. At $\alpha=m^{-2}$, the regularized correction is
\[
 c^{(m)}_n=\frac{m^2}{n(m^2+n^2)},\qquad
 a^\dagger_n-c^{(m)}_n=\frac{n}{m^2+n^2}.
\]
The bias square is exactly the previously computed $B_m$. This links the two cases without conflating their existence hypotheses: the escaping residual of one problem is the regularization bias of an attainable problem.

\begin{corollary}[An explicit noisy attainable repair]\label{WR12}
For the attainable datum $b=(1/n^2)$, suppose $\norm{\widetilde b-b}\le\delta$ with $0<\delta\le1$. Choose
\[
 m=\lceil\delta^{-2/3}\rceil,\qquad\alpha=m^{-2}.
\]
Then the noisy regularized correction satisfies
\[
 \norm{a_\alpha(\widetilde b)-a^\dagger}
 \le\sqrt{\frac2m}+\frac{\delta m}{2}
 \le(\sqrt2+1)\delta^{1/3}.
\]
\end{corollary}
\begin{proof}
The exact-data bias has norm $\sqrt{B_m}\le\sqrt{2/m}$ by WR10. The noise bound in WR09 is $\delta m/2$. The triangle inequality gives the first estimate. The ceiling choice obeys
$\delta^{-2/3}\le m\le\delta^{-2/3}+1\le2\delta^{-2/3}$, which gives the second.
\end{proof}
\written{WR12}

This is a deterministic guarantee for the specified datum and Hilbert noise bound. It is not a statistical confidence statement. The identical parameter rule in the escaping problem cannot produce convergence to an exact protected Hilbert solution, because no such solution exists.

\subsection{What the calculation permits a machine to certify}
The finite program can check a normal equation, exact preservation, and a quadratic optimality identity. It can return a dual witness ruling out a smaller residual at a given budget. In the infinite diagonal example it can return partial sums together with rational tail enclosures. Those are distinct certificates with distinct conclusions.

None of them alone establishes unrestricted semantic realizability. The bounded-budget theorem needs a genuinely uniform bound. The exactification theorem needs control of candidate movement. The regularization limit needs actual attainment of the projected datum. The repeated appearance of a small residual does not replace any of those hypotheses.

The useful next diagnostic is therefore not just ``Did the residual decrease?'' It is ``What happened to the correction norm, what quantity was preserved, and what existence theorem applies to this datum?''
